\documentclass[11pt]{article}

% ===================== PACKAGES =====================
\usepackage[a4paper,margin=1in]{geometry}
\usepackage{amsmath,amssymb}
\usepackage{enumitem}
\usepackage{fancyhdr}
\usepackage{xcolor}

% ===================== HEADER & FOOTER =====================
\pagestyle{fancy}
\fancyhf{}
\rhead{CSE 400: Fundamentals of Probability in Computing}
\lhead{Tutorial 2 -- Complete Scribe}
\rfoot{\thepage}

% ===================== DOCUMENT =====================
\begin{document}

\begin{center}
    \Large \textbf{CSE 400: Fundamentals of Probability in Computing}\\
    \large Tutorial 2 -- Complete Refactored Scribe\\
    \normalsize School of Engineering and Applied Science (SEAS)\\
    February 2026
\end{center}

\vspace{0.5cm}

% ==========================================================
\section{Problem 1: Random Break of a Line Segment}

\subsection*{Modeling and Notation}

Let a point be chosen uniformly on $(0,L)$.  
Let $X$ denote the distance from one endpoint.

\[
X \sim \text{Uniform}(0,L), \quad f_X(x)=\frac{1}{L}
\]

Two segments formed:
\[
X \quad \text{and} \quad L-X
\]

Define ratio:
\[
R = \frac{\min(X,L-X)}{\max(X,L-X)}
\]

\subsection*{Derivation}

We solve:
\[
\frac{\min(X,L-X)}{\max(X,L-X)} < \frac14
\]

Case 1: $X \le L/2$
\[
\frac{X}{L-X} < \frac14
\Rightarrow 4X < L-X
\Rightarrow 5X < L
\Rightarrow X < \frac{L}{5}
\]

Case 2: $X > L/2$
\[
\frac{L-X}{X} < \frac14
\Rightarrow 4L - 4X < X
\Rightarrow 5X > 4L
\Rightarrow X > \frac{4L}{5}
\]

Valid region:
\[
(0, L/5) \cup (4L/5, L)
\]

Total length:
\[
\frac{2L}{5}
\]

\[
P = \frac{2L/5}{L} = \boxed{\frac{2}{5}}
\]

% ==========================================================
\section{Problem 2: Bayes Classification with Normal Densities}

\subsection*{Given}

\[
X|W \sim \mathcal{N}(4,4), \quad X|B \sim \mathcal{N}(6,9)
\]

\[
P(B)=\alpha, \quad P(W)=1-\alpha
\]

Observed: $X=5$

\subsection*{Bayes Theorem}

\[
P(B|X) =
\frac{\alpha f(5|B)}
{\alpha f(5|B)+(1-\alpha)f(5|W)}
\]

Equal error condition:
\[
P(B|5)=\frac12
\]

Using normal density:
\[
f(x)=\frac{1}{\sigma\sqrt{2\pi}}
e^{-\frac{(x-\mu)^2}{2\sigma^2}}
\]

Substitution gives:

\[
\frac{3\alpha}{3\alpha+2(1-\alpha)e^{-5/72}}
=\frac12
\]

Solving:

\[
\boxed{\alpha \approx 0.3827}
\]

% ==========================================================
\section{Problem 3: Optimal Fire Station Location}

\subsection*{(a) Uniform Case}

\[
X \sim \text{Uniform}(0,A)
\]

Minimize:
\[
E|X-a|
\]

After differentiation:

\[
\boxed{a = \frac{A}{2}}
\]

\textbf{Interpretation:} Median minimizes absolute deviation.

\subsection*{(b) Exponential Case}

\[
X \sim \text{Exponential}(\lambda)
\]

\[
\frac{d}{da}E|X-a|=1-2e^{-\lambda a}
\]

\[
e^{-\lambda a}=\frac12
\Rightarrow
\boxed{a=\frac{\ln 2}{\lambda}}
\]

% ==========================================================
\section{Problem 4: Mixture of Exponentials}

Two types:
\[
P(T=i)=p_i
\]

\[
P(X>t)=p_1 e^{-\lambda_1 t}+p_2 e^{-\lambda_2 t}
\]

\[
P(X>t+s|X>t)
=
\frac{
p_1 e^{-\lambda_1(t+s)}+
p_2 e^{-\lambda_2(t+s)}
}{
p_1 e^{-\lambda_1 t}+
p_2 e^{-\lambda_2 t}
}
\]

\textbf{Note:} Mixture destroys memoryless property.

% ==========================================================
\section{Problem 5: Conditional Poisson → Binomial}

\[
X \sim \text{Poisson}(\lambda_1),
\quad
Y \sim \text{Poisson}(\lambda_2)
\]

Find:
\[
P(X=k|X+Y=n)
\]

After simplification:

\[
\boxed{
P(X=k|X+Y=n)
=
\binom{n}{k}
\left(
\frac{\lambda_1}{\lambda_1+\lambda_2}
\right)^k
\left(
\frac{\lambda_2}{\lambda_1+\lambda_2}
\right)^{n-k}
}
\]

\textbf{Result:} Conditional distribution is Binomial.

% ==========================================================
\section{Problem 6: Mixed Distribution Transformation}

\[
X \sim \text{Uniform}[-2,2]
\]

Transformation:

\[
g(x)=
\begin{cases}
x & [-2,-1] \\
0 & (-1,1) \\
x & [1,2]
\end{cases}
\]

\subsection*{PDF}

Continuous parts:

\[
f_Y(y)=\frac14
\]

Point mass:

\[
P(Y=0)=\frac12
\]

\textbf{Important:} Combination of continuous + discrete distribution.

% ==========================================================
\section{Problem 7: Stock Model and CLT Idea}

After 1000 periods:

\[
S_{1000}=S_0 u^k d^{1000-k}
\]

Requirement:
\[
k \ge 470
\]

\[
k \sim \text{Binomial}(1000,0.52)
\]

Probability:
\[
P(k \ge 470)
=
\sum_{k=470}^{1000}
\binom{1000}{k}
(0.52)^k(0.48)^{1000-k}
\]

\textbf{For exam:} Use Normal Approximation.

% ==========================================================
\section{Problem 8: Exponential Repair Time}

\[
f(x)=\frac12 e^{-x/2}
\]

\[
P(X>2)=e^{-1}
\]

Memoryless property:

\[
P(X\ge 10|X>9)=e^{-1/2}
\]

% ==========================================================
\section{Problem 9: Radius Transformation}

\[
Z=\sqrt{X^2+Y^2}
\]

CDF:

\[
F_Z(z)
=
\iint_{x^2+y^2 \le z^2}
f_X(x)f_Y(y)\,dxdy
\]

PDF:

\[
f_Z(z)
=
z
\int_{-z}^{z}
\frac{
f_X(x)\left[
f_Y(\sqrt{z^2-x^2})+
f_Y(-\sqrt{z^2-x^2})
\right]
}{
\sqrt{z^2-x^2}
}
dx
\]

\textbf{Technique:} Differentiate CDF using Leibniz Rule.

% ==========================================================
\section{Problem 10: Joint Distribution of Maxima}

\[
M_n=\max(X_1,\dots,X_n)
\]

\[
P(M_n \le x)=F(x)^n
\]

Joint CDF:

\[
P(M_n \le x, M_{n+1} \le y)
=
\begin{cases}
F(x)^n F(y), & x \le y \\
F(y)^{n+1}, & x > y
\end{cases}
\]

% ==========================================================
\section{Complete Exam Revision Summary}

\begin{itemize}

\item Uniform distributions → geometric probability.
\item Median minimizes $E|X-a|$.
\item Bayes theorem + likelihood ratio.
\item Exponential distribution → memoryless.
\item Mixture ≠ memoryless.
\item Poisson conditional → Binomial.
\item CLT for large Binomial.
\item Transformations may create mixed distributions.
\item Use Leibniz rule for variable-limit integrals.
\item Max of i.i.d. → $F(x)^n$ structure.

\end{itemize}

\bigskip
\centerline{\Large \textbf{End of Tutorial 2 Complete Scribe}}

\end{document}